% Part: incompleteness
% Chapter: arithmetization-syntax
% Section: proofs-in-nd

\documentclass[../../../include/open-logic-section]{subfiles}

\begin{document}

\olfileid[hi]{inc}{art}{pnd}
\olsection{स्वाभाविक निगमन में \usetoken{P}{derivation}}

\begin{explain}
!!{derivation}s का अंकगणितीकरण करने के लिए हमें !!{derivation}s को
संख्याओं के रूप में निरूपित करना होगा। चूँकि !!{derivation}s ऐसे वृक्ष हैं
जिनके नोडों पर !!{formula}s होते हैं और प्रत्येक अनुमान के साथ एक या दो
लेबल जुड़े होते हैं, इसलिए पुनरावर्ती निरूपण सबसे स्पष्ट उपाय है: हम किसी
!!{derivation} को ऐसे टपल के रूप में निरूपित करते हैं जिसके घटक हैं---अंतिम
अनुमान के आधारवाक्यों तक पहुँचने वाले तात्कालिक उप-!!{derivation}s की संख्या,
इन उप-!!{derivation}s के निरूपण, अंतिम !!{formula}, अंतिम अनुमान का
मुक्ति-लेबल, और अंतिम अनुमान का प्रकार बताने वाली एक संख्या।
\end{explain}

\begin{defn}
  यदि $\delta$ स्वाभाविक निगमन में !!a{derivation} है, तो $\Gn{\delta}$
  को निम्न प्रकार आगमनात्मक रूप से परिभाषित किया जाता है:
  \begin{enumerate}
  \item
    यदि $\delta$ में केवल मान्यता~$!A$ है, तो $\Gn{\delta}$,
    $\tuple{0, \Gn{!A}, n}$ है। यदि यह !!{undischarged} मान्यता है तो
    संख्या~$n$, $0$ है; अन्यथा वह उसका संख्यात्मक लेबल है।
  \item
    यदि $\delta$ का अंत शून्य, एक, दो अथवा तीन आधारवाक्यों वाले अनुमान में
    होता है, तो $\Gn{\delta}$ क्रमशः
    \begin{align*}
      & \tuple{0, \Gn{!A}, n, k}, \\
      & \tuple{1, \Gn{\delta_1}, \Gn{!A}, n, k}, \\
      & \tuple{2, \Gn{\delta_1}, \Gn{\delta_2}, \Gn{!A}, n, k}, \text{ अथवा}\\
      & \tuple{3, \Gn{\delta_1}, \Gn{\delta_2}, \Gn{\delta_3}, \Gn{!A}, n,
        k}
    \end{align*}
    है। यहाँ $\delta_1$, $\delta_2$, $\delta_3$ वे उप-!!{derivation}s हैं
    जिनका अंत $\delta$ के अंतिम अनुमान के आधारवाक्य (अथवा आधारवाक्यों) में
    होता है; $!A$, $\delta$ के अंतिम अनुमान का निष्कर्ष है; $n$ अंतिम
    अनुमान का मुक्ति-लेबल है (यदि अनुमान किसी मान्यता को मुक्त नहीं करता,
    तो $0$); और अंतिम अनुमान में प्रयुक्त नियम के अनुसार $k$ निम्न तालिका
    से मिलता है।

    \begin{tabular}{lccccccc}
      \text{नियम:} & \Intro{\land} & \Elim{\land} & \Intro{\lor} & \Elim{\lor} \\
      $k$: & 1 & 2 & 3 & 4  \\[1.5ex]
      \text{नियम:} &  \Intro{\lif} & \Elim{\lif} & \Intro{\lnot} & \Elim{\lnot} \\
      $k$: & 5 & 6 & 7 & 8 \\[1.5ex]
      \text{नियम:}  &  \FalseInt & \FalseCl & \Intro{\lforall} & \Elim{\lforall} \\
      $k$: & 9 & 10 & 11 & 12 \\[1.5ex]
      \text{नियम:} & \Intro{\lexists} & \Elim{\lexists} & \Intro{\eq} & \Elim{\eq} \\
      $k$: & 13 & 14 & 15 & 16
    \end{tabular}
  \end{enumerate}
\end{defn}

\begin{ex}
  इस अत्यंत सरल !!{derivation} पर विचार करें:
  \begin{prooftree}
    \AxiomC{$\Discharge{!A \land !B}{1}$}
    \RightLabel{\Elim{\land}}
    \UnaryInfC{$!A$}
    \DischargeRule{\Intro{\lif}}{1}
    \UnaryInfC{$(!A \land !B) \lif !A$}
  \end{prooftree}
  मान्यता की ग्योडेल (G\"odel) संख्या $d_0 = \tuple{0,
    \Gn{!A \land !B}, 1}$ होगी। $\Elim{\land}$ के निष्कर्ष पर समाप्त होने
  वाली !!{derivation} की ग्योडेल (G\"odel) संख्या $d_1 = \tuple{1,
     d_0, \Gn{!A}, 0, 2}$ होगी ($1$, क्योंकि $\Elim{\land}$ का एक
  आधारवाक्य है; निष्कर्ष~$!A$ की ग्योडेल (G\"odel) संख्या; $0$, क्योंकि
  कोई मान्यता मुक्त नहीं होती; और $2$, $\Elim{\land}$ को कूटबद्ध करने वाली
  संख्या है)। तब पूरी !!{derivation} की ग्योडेल (G\"odel) संख्या
  $\tuple{1, d_1, \Gn{((!A \land !B) \lif
      !A)}, 1, 5}$ है, अर्थात्
  \[
  \tuple{1, \tuple{1, \tuple{0, \Gn{(!A \land !B)}, 1}, \Gn{!A}, 0, 2},
    \Gn{((!A \land !B) \lif !A)}, 1, 5}.
  \]
\end{ex}

\begin{explain}
!!{derivation}s का एक निरूपण चुन लेने के बाद हमें यह भी दिखाना होगा कि
ऐसी !!{derivation}s की ग्योडेल (G\"odel) संख्याओं को हम आदिम पुनरावर्ती
ढंग से परिचालित कर सकते हैं और उनके आवश्यक गुण तथा संबंध व्यक्त कर सकते
हैं। कुछ संक्रियाएँ सरल हैं: उदाहरणतः, !!a{derivation} की ग्योडेल (G\"odel)
संख्या~$d$ दिए जाने पर $\fn{EndFmla}(d) = (d)_{(d)_0+1}$ उसके अंतिम
!!{formula} की ग्योडेल (G\"odel) संख्या देता है,
$\fn{DischargeLabel}(d) = (d)_{(d)_0 + 2}$ उसका मुक्ति-लेबल देता है और
$\fn{LastRule}(d) = (d)_{(d)_0 + 3}$ अंतिम अनुमान का प्रकार बताने वाली
संख्या देता है। कुछ संक्रियाएँ कहीं कठिन हैं। हम कम-से-कम यह रेखांकित
करेंगे कि उन्हें कैसे किया जा सकता है। लक्ष्य यह दिखाना है कि संबंध
``$\delta$,~$\Gamma$ से~$!A$ की !!a{derivation} है'' $\delta$ और~$!A$ की
ग्योडेल (G\"odel) संख्याओं का आदिम पुनरावर्ती संबंध है।
\end{explain}

\begin{prop}
  निम्न संबंध आदिम पुनरावर्ती हैं:
  \begin{enumerate}
  \item $!A$,~$\delta$ में लेबल~$n$ वाली मान्यता के रूप में आता है।
  \item $\delta$ में लेबल~$n$ वाली सभी मान्यताएँ~$!A$ के रूप की हैं
    (अर्थात् हम $\delta$ में लेबल~$n$ का उपयोग करके मान्यता~$!A$ को
    !!{discharge} कर सकते हैं)।
  \end{enumerate}
\end{prop}

\begin{proof}
  हमें यह दिखाना है कि !!{formula}s की ग्योडेल (G\"odel) संख्याओं और
  !!{derivation}s की ग्योडेल (G\"odel) संख्याओं के बीच संगत संबंध आदिम
  पुनरावर्ती हैं।
  \begin{enumerate}
  \item हम दिखाना चाहते हैं कि $\fn{Assum}(x, d, n)$ आदिम पुनरावर्ती है;
    यह संबंध तभी लागू होता है जब $x$, ग्योडेल (G\"odel) संख्या~$d$ वाली
    !!{derivation} की, लेबल~$n$ वाली मान्यता की ग्योडेल (G\"odel) संख्या
    हो। ऐसा तभी है जब ग्योडेल (G\"odel) संख्या~$\tuple{0, x, n}$ वाली
    !!{derivation},~$d$ की उप-!!{derivation} हो। ध्यान दें कि !!{derivation}s
    का हमारा कूटलेखन
    \olref[cmp][rec][tre]{sec} में दिए गए वृक्ष-कूटलेखन का विशेष मामला है,
    इसलिए आदिम पुनरावर्ती फलन $\fn{SubtreeSeq}(d)$,~$d$ की सभी
    उप-!!{derivation}s की ग्योडेल (G\"odel) संख्याओं का एक अनुक्रम देता है
    (जिसकी लंबाई अधिकतम~$d$ है)। अतः हम परिभाषित कर सकते हैं
    \[
    \fn{Assum}(x, d, n) \defiff \bexists{i<d}{(\fn{SubtreeSeq}(d))_i =
      \tuple{0, x, n}}.
    \]
  \item हम दिखाना चाहते हैं कि $\fn{Discharge}(x, d, n)$ आदिम पुनरावर्ती
    है; यह संबंध तभी लागू होता है जब ग्योडेल (G\"odel) संख्या~$d$ वाली
    !!{derivation} में लेबल~$n$ वाली सभी मान्यताएँ, ग्योडेल (G\"odel)
    संख्या~$x$ वाले !!{formula} की ही हों। किंतु यह संबंध तभी और केवल तभी
    लागू होता है जब $\bforall{y<d}{(\fn{Assum}(y, d, n) \lif y = x)}$।
  \end{enumerate}
\end{proof}

\begin{prop}
  \ollabel{prop:followsby}
  वह गुण $\fn{Correct}(d)$ आदिम पुनरावर्ती है, जो तभी लागू होता है जब
  ग्योडेल (G\"odel) संख्या~$d$ वाली !!{derivation}~$\delta$ का अंतिम
  अनुमान सही हो।
\end{prop}

\begin{proof}
  यहाँ हमें प्रत्येक अनुमान-नियम~$R$ के लिए यह दिखाना है कि संबंध
  $\fn{FollowsBy}_R(d)$ आदिम पुनरावर्ती है, जहाँ
  $\fn{FollowsBy}_R(d)$ तभी लागू होता है जब $d$, !!{derivation}~$\delta$ की
  ग्योडेल (G\"odel) संख्या हो और $\delta$ का अंतिम !!{formula},~$\delta$
  की तात्कालिक उप-!!{derivation}s से~$R$ के सही अनुप्रयोग द्वारा निकले।

  \Intro{\land} नियम का मामला सरल है। यदि $\delta$ का अंत सही
  $\Intro{\land}$ अनुमान में होता है, तो वह ऐसा दिखता है:
  \begin{prooftree}
    \AxiomC{}
    \RightLabel{$\delta_1$}
    \DeduceC{$!A$}

    \AxiomC{}
    \RightLabel{$\delta_2$}
    \DeduceC{$!B$}

    \RightLabel{\Intro\land}
    \BinaryInfC{$!A \land !B$}
  \end{prooftree}
  तब $\delta$ की ग्योडेल (G\"odel) संख्या~$d$, $\tuple{2, d_1, d_2,
    \Gn{(!A \land !B)}, 0, k}$ है, जहाँ $\fn{EndFmla}(d_1) = \Gn{!A}$,
  $\fn{EndFmla}(d_2) = \Gn{!B}$, $n=0$, और $k=1$। अतः हम
  $\fn{FollowsBy}_{\Intro\land}(d)$ को इस प्रकार परिभाषित कर सकते हैं:
  \begin{multline*}
    (d)_0 = 2 \land \fn{DischargeLabel}(d) = 0 \land \fn{LastRule}(d) = 1 \land {}\\
    \fn{EndFmla}(d) = {}\\
    \Gn{(} \concat \fn{EndFmla}((d)_1) \concat \Gn{\land}
    \concat \fn{EndFmla}((d)_2) \concat \Gn{)}.
  \end{multline*}

  दूसरा सरल उदाहरण $\Intro\eq$ नियम है। इसका कोई आधारवाक्य नहीं है,
  इसलिए मान्यताओं की तरह $(d)_0 = 0$। इसका कोई मुक्ति-लेबल भी नहीं है,
  अर्थात् $n=0$। किंतु किसी बंद पद~$t$ के लिए $!A$,~$\eq[t][t]$ के रूप
  का होना चाहिए। यहाँ एक आदिम पुनरावर्ती परिभाषा यह है:
  \begin{multline*}
    (d)_0 = 0 \land \fn{DischargeLabel}(d) = 0 \land {}\\
    \bexists{t<d}{(\fn{ClTerm}(t) \land
      \fn{EndFmla}(d) = {}\\
      \Gn{{\eq}(} \concat t \concat \Gn{,} \concat t \concat \Gn{)})}.
  \end{multline*}

  एक अधिक जटिल उदाहरण में, $\fn{FollowsBy}_{\Intro{\lif}}(d)$ तभी लागू
  होता है जब $\delta$ का अंतिम !!{formula},~$(!A \lif !B)$ के रूप का हो,
  $\delta_1$ का अंतिम !!{formula},~$!B$ हो और $\delta$ में लेबल~$n$ वाली
  हर मान्यता~$!A$ के रूप की हो। इसे हम निम्न प्रकार आदिम पुनरावर्ती रूप
  में व्यक्त कर सकते हैं:
  \begin{multline*}
    (d)_0 = 1 \land {}\\
    \bexists{a<d}{(\fn{Discharge}(a, (d)_1, \fn{DischargeLabel}(d)) \land {}}\\
      \fn{EndFmla}(d) = (\Gn{(} \concat a \concat \Gn{\lif}
      \concat \fn{EndFmla}((d)_1) \concat \Gn{)}))
  \end{multline*}
  ($a$ को~$!A$ की ग्योडेल (G\"odel) संख्या मानें)।

  एक और उदाहरण के लिए \Intro{\lexists} पर विचार करें। यहाँ $\delta$ का
  अंतिम अनुमान तभी सही है जब कोई !!a{formula}~$!A$, कोई बंद पद~$t$ और कोई
  !!a{variable}~$x$ ऐसे हों कि $\Subst{!A}{t}{x}$,
  !!{derivation}~$\delta_1$ का अंतिम !!{formula} हो और
  $\lexists[x][!A]$ अंतिम अनुमान का निष्कर्ष हो। अतः
  $\fn{FollowsBy}_{\Intro{\lexists}}(d)$ तभी लागू होता है जब
  \begin{multline*}
    (d)_0 = 1 \land \fn{DischargeLabel}(d) = 0 \land {} \\
    \bexists{a < d}{\bexists{x<d}{\bexists{t<d}{
          (\fn{ClTerm}(t) \land \fn{Var}(x)  \land {}}}}\\
    \fn{Subst}(a,t,x) = \fn{EndFmla}((d)_1) \land
    \fn{EndFmla}(d) = (\Gn{\lexists} \concat x \concat a)).
  \end{multline*}

  तब हम $\fn{Correct}(d)$ को इस प्रकार परिभाषित करते हैं:
  \begin{multline*}
    \fn{Sent}(\fn{EndFmla}(d)) \land {}\\
    (\fn{LastRule}(d) = 1 \land
    \fn{FollowsBy}_{\Intro\land}(d)) \lor \dots \lor {}\\
    (\fn{LastRule}(d) = 16 \land \fn{FollowsBy}_{\Elim\eq}(d)) \lor {}\\
    \bexists{n<d}{\bexists{x<d}{(d = \tuple{0, x, n})}}.
  \end{multline*}
  पहली पंक्ति सुनिश्चित करती है कि~$d$ का अंतिम !!{formula} कोई वाक्य है।
  अंतिम पंक्ति उस मामले को समेटती है जिसमें $d$ केवल एक मान्यता है।
\end{proof}

\begin{prob}
  निम्न गुणों को \olref[inc][art][pnd]{prop:followsby} की भाँति परिभाषित
  करें:
  \begin{enumerate}
  \item $\fn{FollowsBy}_{\Elim{\lif}}(d)$,
  \item $\fn{FollowsBy}_{\Elim{\eq}}(d)$,
  \item $\fn{FollowsBy}_{\Elim{\lor}}(d)$,
  \item $\fn{FollowsBy}_{\Intro{\lforall}}(d)$।
  \end{enumerate}
  अंतिम गुण के लिए आपको यह भी दिखाना होगा कि हम आदिम पुनरावर्ती ढंग से
  जाँच सकते हैं कि ग्योडेल (G\"odel) संख्या~$d$ वाली !!{derivation} का अंतिम
  अनुमान स्वचर-प्रतिबंध को संतुष्ट करता है या नहीं; अर्थात्
  $\Intro{\lforall}$ अनुमान का स्वचर~$a$, न तो~$d$ के अंतिम !!{formula} में
  और न~$d$ की किसी खुली मान्यता में आता है। इसके लिए आप
  \olref[inc][art][pnd]{prop:openassum} का आदिम पुनरावर्ती विधेय
  $\fn{OpenAssum}$ उपयोग कर सकते हैं।
\end{prob}

\begin{prop}
  \ollabel{prop:deriv}
  वह संबंध $\fn{Deriv}(d)$ आदिम पुनरावर्ती है, जो तब लागू होता है जब $d$
  किसी सही !!{derivation}~$\delta$ की ग्योडेल (G\"odel) संख्या हो।
\end{prop}

\begin{proof}
  !!^a{derivation}~$\delta$ तभी सही है जब उसका प्रत्येक अनुमान किसी नियम
  का सही अनुप्रयोग हो, अर्थात् जब उसकी प्रत्येक उप-!!{derivation} का अंत
  किसी सही अनुमान में हो। अतः $\fn{Deriv}(d)$ तभी लागू होता है जब
  \[
  \bforall{i<\len{\fn{SubtreeSeq}(d)}}{\fn{Correct}((\fn{SubtreeSeq}(d))_i)}
  \]
\end{proof}

\begin{prop}
  \ollabel{prop:openassum} संबंध $\fn{OpenAssum}(z, d)$ तब लागू होता है जब
  $z$, ग्योडेल (G\"odel) संख्या~$d$ वाली !!{derivation}~$\delta$ की किसी
  !!{undischarged} मान्यता~$!A$ की ग्योडेल (G\"odel) संख्या हो; यह संबंध
  आदिम पुनरावर्ती है।
\end{prop}

\begin{proof}
  मान्यता की कोई उपस्थिति तभी !!{discharged} होती है जब वह~$\delta$ की ऐसी
  उप-!!{derivation} में लेबल~$n$ के साथ आती है, जिसका अंत मुक्ति-लेबल~$n$
  वाले नियम में होता है। अतः $!A$,~$\delta$ की !!a{undischarged} मान्यता
  तभी है जब उसकी कम-से-कम एक उपस्थिति~$\delta$ में !!{discharged} न हो।
  हमें सावधान रहना चाहिए: $\delta$ में~$!A$ की !!{discharged} और
  !!{undischarged}, दोनों प्रकार की उपस्थितियाँ हो सकती हैं।

  ऐसे अनुक्रम $\delta_0$, \dots, $\delta_k$ पर विचार करें जिसमें
  $\delta_0 = \delta$, $\delta_k$ मान्यता $\Discharge{!A}{n}$ है
  (किसी~$n$ के लिए), और $\delta_{i+1}$,~$\delta_i$ की तात्कालिक
  उप-!!{derivation} है। यदि ऐसा अनुक्रम मौजूद है जिसमें किसी~$\delta_i$
  का अंत मुक्ति-लेबल~$n$ वाले अनुमान में नहीं होता, तो $!A$,~$\delta$ की
  !!a{undischarged} मान्यता है।

  आदिम पुनरावर्ती फलन $\fn{SubtreeSeq}(d)$ हमें~$\delta$ की सभी
  उप-!!{derivation}s की ग्योडेल (G\"odel) संख्याओं का अनुक्रम देता है।
  $\delta$ की उप-!!{derivation}s की ग्योडेल (G\"odel) संख्याओं का प्रत्येक
  अनुक्रम इसका उपानुक्रम है। उपानुक्रम होना आदिम पुनरावर्ती संबंध है:
  $\fn{Subseq}(s, s')$ तभी लागू होता है जब
  $\bforall{i<\len{s}}{\lexists[j<\len{s'}][(s)_i = (s')_j]}$। तात्कालिक
  उप-!!{derivation} होना भी ऐसा ही है: $\fn{Subderiv}(d, d')$ तभी लागू होता
  है जब $\bexists{j<(d')_0}{d = (d')_j}$। अतः हम
  $\fn{OpenAssum}(z, d)$ को इस प्रकार परिभाषित कर सकते हैं:
  \begin{multline*}
    \bexists{s<\fn{SubtreeSeq}(d)}{(\fn{Subseq}(s, \fn{SubtreeSeq}(d))
      \land (s)_0 = d \land {}} \\
    \bexists{n<d}{((s)_{\len{s} \tsub 1} = \tuple{0, z, n} \land {}}\\
      \bforall{i<(\len{s} \tsub 1)}{(\fn{Subderiv}((s)_{i+1}, (s)_i) \land {}}\\
      \fn{DischargeLabel}((s)_i) \neq n))).
  \end{multline*}
\end{proof}

\begin{prop}
  \ollabel{prop:prf-prim-rec}
  मान लें कि $\Gamma$, !!{sentence}s का आदिम पुनरावर्ती समुच्चय है। तब
  संबंध $\Prf[\Gamma](x, y)$, जो व्यक्त करता है कि ``$x$,~$\Gamma$ में
  स्थित !!{undischarged} मान्यताओं से~$!A$ की !!a{derivation}~$\delta$ का
  कूट है और $y$,~$!A$ की ग्योडेल (G\"odel) संख्या है'', आदिम पुनरावर्ती है।
\end{prop}

\begin{proof}
  मान लें कि ``$y \in \Gamma$'' आदिम पुनरावर्ती विधेय~$R_\Gamma(y)$ द्वारा
  दिया गया है। हमें दिखाना है कि $\Prf[\Gamma](x, y)$ आदिम पुनरावर्ती है;
  यह संबंध तभी लागू होता है जब $y$ किसी वाक्य~$!A$ की ग्योडेल (G\"odel)
  संख्या हो और $x$, ऐसे स्वाभाविक निगमन !!{derivation} का कूट हो जिसका
  अंतिम !!{formula}~$!A$ है तथा जिसकी सभी !!{undischarged} मान्यताएँ
  $\Gamma$ में हैं।

  \olref{prop:deriv} के अनुसार वह गुण $\fn{Deriv}(x)$ आदिम पुनरावर्ती है,
  जो तभी लागू होता है जब $x$ स्वाभाविक निगमन की किसी सही
  !!{derivation}~$\delta$ की ग्योडेल (G\"odel) संख्या हो। अतः हम
  $\Prf[\Gamma](x, y)$ को इस प्रकार परिभाषित कर सकते हैं:
  \begin{align*}
    \Prf[\Gamma](x, y) \defiff {}
    & \fn{Deriv}(x) \land \fn{EndFmla}(x) = y \land {} \\
    & \bforall{z < x}{(\fn{OpenAssum}(z, x) \lif R_\Gamma(z))}.
  \end{align*}
\end{proof}

\end{document}
